% =====================================================================
%  EXPERIMENTELLES  (Abschnitt 1.3 der Anleitung)
%  Skizze des Aufbaus + kurze, aber genaue Schilderung der Durchführung.
%  Alle vorgenommenen Einstellungen protokollieren!
% =====================================================================
\section{Experimentelles}
\label{sec:experimentelles}

\subsection{Versuchsaufbau}
Beschreibung und Skizze/Foto des experimentellen Aufbaus. Fremde
Abbildungen sind mit Quellenangabe zu kennzeichnen (Abschnitt 3 der
Anleitung).

\begin{figure}[H]
  \centering
  % \includegraphics[width=0.8\textwidth]{aufbau.pdf}
  \fbox{\parbox[c][5cm][c]{0.7\textwidth}{\centering Platzhalter: Skizze/Foto des Versuchsaufbaus}}
  \caption{Versuchsaufbau mit allen relevanten Komponenten.}
  \label{fig:aufbau}
\end{figure}

Verwendete Geräte und Materialien inkl. Inventarnummern und
Einstellungen (z. B. Verstärkung, Zeitbase, Messbereiche):

\begin{table}[H]
  \centering
  \caption{Verwendete Geräte und wichtige Einstellungen.}
  \label{tab:geraete}
  \begin{tabular}{llll}
    \toprule
    Gerät & Typ/Hersteller & Inventarnummer & Einstellung \\
    \midrule
    Oszilloskop & Mustertyp A & 12345 & 2\,ms/div \\
    Multimeter & Mustertyp B & 67890 & Bereich 10\,V \\
    Spannungsquelle & Mustertyp C & 11121 & 5,0\,V \\
    \bottomrule
  \end{tabular}
\end{table}

\subsection{Durchführung}
Kurze, aber genaue Schilderung des Ablaufs:
\begin{itemize}
  \item Vorbereitung und Justage (z. B. Nullabgleich, Kalibrierung).
  \item Messung: Was wurde wie und wie oft gemessen (N Wiederholungen)?
  \item Alle vorgenommenen Einstellungen und Änderungen notieren.
\end{itemize}

Hinweis zur Schreibweise (Abschnitt 2 der Anleitung): neutrale
Erzählperspektive (keine Ich-Form), kurz und prägnant, keine Redundanzen.
